% !TEX root = main2.tex
\section{Market-data and calibration pipeline}
\label{app:market-pipeline}

This appendix records the data and numerical pipeline behind
Section~\ref{sec:estimation}. It follows the broader SPX pipeline described in
\cite{DYSANOS2026}, restricted here to the short-dated fixed-expiry SANOS
experiment. The complete transformation is
\[
\begin{aligned}
\text{raw IvyDB quotes}
&\longrightarrow \text{daily cleaning and parity-forward enrichment}\\
&\longrightarrow \text{arbitrage-free SANOS eligibility surface}\\
&\longrightarrow \text{European short-horizon quote selection}\\
&\longrightarrow \text{daily 25-by-200 calibration tensor}\\
&\longrightarrow \text{direct or affine end-to-end fit}.
\end{aligned}
\]

\subsection{Raw data and cleaning}
\label{app:pipeline-raw}

The source is OptionMetrics IvyDB US data accessed through WRDS. For each
eligible SPX market date we load the index level and previous index level,
option strikes and calendar expiries, call and put bid and ask quotes, volume,
open interest, AM/PM settlement flags, and the zero curve. The retained sample
runs from 3 January 2022 through 29 August 2025 and contains 891 valid full-day
NYSE sessions.

The raw cleaning step removes known bad ticker--date observations, missing or
invalid quotes, invalid settlement flags, same-day expiries, irregular
expiries, and quotes whose bid does not exceed $0.011$. Business days and year
fractions to expiry are then computed; an AM-settled expiry is treated as
$0.2$ of a business day earlier than its corresponding PM expiry.

For each expiry we recompute the forward from put--call parity rather than use
the delayed IvyDB forward table. With discount factor $D$, define the
discounted-forward bounds at a strike $K$ by
\[
  G^{\mathrm{ask}}(K)=C^{\mathrm{ask}}(K)-P^{\mathrm{bid}}(K)+DK,
  \qquad
  G^{\mathrm{bid}}(K)=C^{\mathrm{bid}}(K)-P^{\mathrm{ask}}(K)+DK.
\]
Their midpoint is averaged across strikes with inverse-spread weights,
\[
  \widehat G
  =\frac{\sum_K w_K\,[G^{\mathrm{ask}}(K)+G^{\mathrm{bid}}(K)]/2}
         {\sum_K w_K},
  \qquad
  w_K=\frac1{\max\{G^{\mathrm{ask}}(K)-G^{\mathrm{bid}}(K),10^{-3}\}}.
\]
The resulting forward is used to construct the tightest call-equivalent quote,
forward-normalized strikes and prices, implied volatilities, and Greeks. An
expiry must contain at least ten strikes with a traded call or put and at least
four strikes with both sides traded. The normalized spread must exceed
$10^{-6}$ and retained implied volatilities must lie between $0.01$ and $5$.

\subsection{Static preprocessing and quote selection}
\label{app:pipeline-static}

Before applying the model-specific horizon, the pipeline retains quotes with
vega per unit total volatility of at least $10^{-3}$, log volume and log open
interest of at least one, and at least ten surviving options per expiry. The
nominal per-day quote cap is 100,000 and is therefore non-binding here.

For data-quality control and to construct a stable ATM curve, every daily book
is passed through the linear SANOS preprocessing of \cite{SANOS2026}. The
strike grid is extended, each expiry is fitted inside the quoted bid--ask
intervals, negligible-density redundant strikes are removed, and the joint
surface fit enforces both butterfly and calendar-spread constraints. The
smoothness parameter is $0.25$. Expiries and strikes for which this
arbitrage-free auxiliary fit is unavailable are not passed to the final
calibration.

This auxiliary SANOS surface determines eligibility, the ATM term structure,
and the floating model support. It is not the calibration target: the JAX fits
in Section~\ref{sec:estimation-fit} use the original retained call and put
bid--ask intervals. No linear-programming state or fitted midpoint is supplied
as a target.

The final selection keeps European quotes no farther than 25 business days
from the observation date, excludes the $0.2$-business-day slice, and retains
the standardized log-moneyness interval $[-2.5,1.5]$. This leaves 2,192,458
quotes, between 1,590 and 5,096 per date.

\subsection{Daily calibration tensor}
\label{app:pipeline-tensor}

The model expiry grid is
\[
  T_j=\frac{j}{252},\qquad j=1,\ldots,25.
\]
For each date, 100 interior-strike candidates are generated from the front
ATM scale and 100 from the 25-day ATM scale. They are merged with a minimum
log-spacing safeguard to form one 200-node support shared across all model
expiries on that date. The support floats from date to date. Two boundary
nodes, equal to one half of the lowest inner strike and one and a half times
the highest inner strike, complete the 202-node density grid.

Each retained market quote stores its call/put flag, normalized strike, bid,
ask, midpoint, business days to expiry, adjacent model-expiry indices, and
linear interpolation weight. The contemporaneous SPX level is retained as a
reference spot. The immutable calibration tensor pads each date to a multiple
of 128 quotes and carries a Boolean mask, so padding never contributes to the
objective.

At valuation time the DLV recursion produces the 25 discrete densities. The
fixed 5\% Black factor convolves each discrete density into a smooth call-price
curve, and the two adjacent expiry curves are linearly interpolated to the
market quote's exact expiry. Puts follow from put--call parity.

\subsection{Optimization and retained artifacts}
\label{app:pipeline-optimization}

The direct and affine parameterizations are those in
Section~\ref{sec:estimation-fit}. Direct states are initialized from the
date-specific ATM term structure and are constant across strike at
initialization. Affine means use the average of these initial states; codes and
loadings receive small centered Gaussian initial values. DLV squared is bounded
between $10^{-4}$ and $9$, with a differentiable barrier near the upper bound.

The loss in \eqref{eq:estimation-quote-objective} is differentiated over every
retained quote. Quote blocks of size 128 are used only to limit accelerator
memory: their gradients are summed before each Adam update, so this is an exact
full-sample gradient rather than stochastic minibatching. Each fit uses 300
initial epochs with a cosine learning-rate schedule from $2\times10^{-2}$ to
$10^{-3}$, followed by a 700-epoch continuation from $5\times10^{-3}$ to
$5\times10^{-5}$. Checkpoints store either all daily direct states or the
affine mean, loadings, and daily codes. The retained reports recompute every
quote from the final checkpoint and produce the statistics in
Table~\ref{tab:numerics-static-market-fit}.
